Настоящая глава представляет собой обзор существующих решений для извлечения библиографической информации, структурированной по нескольким направлениям: препроцессинг и детекция разделов, сегментация и парсинг ссылок, а также применение больших языковых моделей. Отдельно обсуждаются проблемы, связанные с неструктурированностью PDF, стилевым разнообразием оформления и недостаточно проработкой русскоязычных корпусов. На основе анализа современных подходов формулируются требования к разрабатываемой системе и обосновывается целесообразность использования LoRA-добучения.

\section{Проблемы извлечения библиографии}

При решении задачи извлечения библиографии необходимо решать следующие задачи\cite{system_higher}:
\begin{itemize}
    \item Проблема «грязных» данных: OCR-ошибки, артефакты парсинга, повреждённые PDF.
    \item Стилевой хаос: ГОСТ, amsbib, APA, MLA, Vancouver, Chicago…
    \item Спецсимволы, как греческие буквы, формулы в заголовках, разделение авторов различной пункт
    \item Формулы в библиографических записях
    \item Где в документе заканчивается текст и начинаются ссылки? 
    \item Где заканчивается одна ссылка и начинается другая?     
    \item Неоднозначности внутри полей: инициалы авторов, фамилиии с частицами, разделение авторов различной пунктуацией или союзами.
    \item Применение больших языковых моделей несёт риск галлюцинаций.
\end{itemize}

Отдельно стоит упомянуть проблему извлечения текстов из pdf-файлов. Эта задача не рассматривается отдельно в данной работе, но она довольно важна в контексте общего конвейера. 

\subsection{Проблемы извлечения текста}
Прежде всего для последующего анализа научной публикации необходимо извлечь текст из статьи. В современной практики научные публикации хранятся в формате Portable Document Format (pdf), что располагает к совокупности возможных ошибок. Ключевая особенность формата заключается в том, что формат PDF разработан для визуальной фиксации, а не для анализа информации в нём имеющейся\cite{bast2017benchmark}. Из-за этой особенности возникает разрыв между тем, что видит читатель и тем, что получает программный парсер. С другой стороны именно за счёт этого противоречия PDF-файлы завоевали свою популярность, так как сохраняется фиксированная пагинация в виде нумерации страниц, а также гарантия, что структурные единицы, как рисунки, формы, таблицы, -- одинаково отображаются почти на всех устройствах. Проблему отсутствия структурной информации решают специальные форматы, например, tagged pdf, но научные публикации, как правило, распространяются без такой разметки. 

Одним из источников проблем является неоднородность структуры pdf-файлов. В зависимости от системы генерации документ может содержать текстовый слой, например, при экспорте из latex или Microsoft Word, или представлять собой растровые изображения копий страниц, полученные с помощью сканирования, или являться промежуточной стадией, когда текстовый слой есть, но логической структуры нет. Если pdf-файл получен исключительно за счёт сканирования, то тут уже никак не обойтись без применения методов оптического распознавания символов (OCR). Даже файлы с текстовым слоем могут содержать символы в нестандартных кодировках или иметь встроенные собственные шрифты без соответствия их с Unicode -- всё это приводит к артефактам парсинга и получению нечитаемых последовательностей. Фундаментальной причиной символьных искажений является архитектура шрифтовых подсистем PDF. При создании документа издательские системы часто применяют включение в файл только использованных глифов и генерируют собственные таблицы ToUnicode CMap сопоставления кодов символов. Если такая таблица отсутствует или сформирована некорректно, парсер получает последовательность байтов без семантической привязки к Unicode. Особенно критично это для кириллицы, где разные генераторы: LaTeX с T2A, Word с Windows-1251, старые верстальные системы, -- используют несовместимые кодировки. В результате фамилии авторов, названия журналов и годы издания преобразуются в нечитаемые последовательности, что делает невозможным их последующее сопоставление с библиометрическими базами данных. Причём некоторые pdf-файлы могут быть закодированы, а часть храниться в сети Интернет в повреждённом виде, что сводит к нулю возможность извлечения качественной информации.

PDF не является семантическим форматом, в отличие от html или xml, где текст может быть размечен по смыслу: заголовки, абзацы, таблицы. При программном извлечении информации с помощью библиотек-парсеров: PDFplumber, PyMuPDF, PDFBox и д.р\cite{hong2021challenges}, -- могут возникать некоторые ошибки. Популярные артефакты: разрыв слов по переносам между строками, нарушения порядка предложений при двухколонной вёрстке, потеря математических символов или других специальных символов, как греческие буквы, текстовые блоки без семантической последовательности. Например, фамилия <<Красоткин>> при парсинге может быть извлечена, как два слова <<Красот>> на одной строке и <<кин>> на следующей строке. При последующем выделение библиографических областей это создаст ложное впечатление о границе между токенами. Особую сложность представляет восстановление логического порядка чтения в многостолбцовой вёрстке. Стандартные библиотеки экстракции часто объединяют текстовые блоки по координатам Y, игнорируя визуальные колонки, что приводит к перемешиванию соседних библиографических записей. Дополнительно проблему усложняют плавающие элементы: переносы ссылок между страницами, маркировка продолжения (продолжение..., cont.), а также вложенные структуры, когда внутри одной записи встречаются цитаты на цитаты. Без этапа layout-анализа такие артефакты интерпретируются как разрывы записей или слияние нескольких источников в один токен.

Старые выпуски журналов нередко распространяются в виде сканированных документов, поэтому для работы с ними единственный выход это использование OCR\cite{sumathi2012techniques}. Современные движки, как EasyOCR, Tesseract или коммерческое решение ABBYY FineReader, демонстрируют высокое качество распознавания на чётких сканах с горизонтальным расположением строк. Однако, когда в документах появляются шумы, сканы низкого качества, даже меньше 200 dpi, нестандартная ориентация страниц, используются редкие шрифты с засечками, то их эффективность резко падает. Вызовом для OCR при работе с русским языком являются буквы: <<ё>>, <<ь>>, <<ъ>>, -- а также кириллические сокращения, что может быть воспринято как латиница.   

При извлечение текста из pdf-файлов стоит понимать, что ошибки на этом этапе несут кумулятивный характер\cite{tkaczyk2015cermine}. Если не удаётся качественно выделить текстовый слой, то тогда не получится найти границы библиографии; если возникают артефакты при извлечение текста при переходе между строками, то может потеряться библиографическая запись; если неверно распознается фамилия автора в списке литературы, то полученный результат не получится сопоставить с библиометрическими базами данных. Ошибки препроцессинга носят строго каскадный характер и нелинейно усиливаются на последующих этапах. Даже незначительные искажения на уровне символов: замена 0 на O, пропуск дефиса в составной фамилии, склеивание инициалов, -- приводят к отказу систем нормализации: CrossRef API, DOI resolver, ORCID matching. В результате формируются «сиротские» узлы в библиометрических графах, искажаются метрики соавторства и цитируемости, а автоматические системы рекомендации теряют точность. Это обуславливает необходимость либо внедрения устойчивых к шуму алгоритмов посткоррекции, либо использования моделей, способных восстанавливать семантику на уровне записи, а не на уровне отдельных символов.

\subsection{Проблемы детекции позиции начала и конца библиографии}

После получения текстового слоя из документа следующий вопрос заключается в том, чтобы обнаружить, где в документе заканчивается основной текст и начинаются ссылки. На первый взгляд, кажется, что достаточно найти заголовок списка литературы, задействуя маски, например <<Список литературы>>, References, <<Библиографический список>>. Тем не менее на практике задача усугубляется несколькими факторами.

После парсинга документа в полученном тексте может присутствовать шум. Такой как: номеры страниц, колонтитулы, технические метрики. Так если на этом этапе позиция библиографии ищется по строгим правилам, например References, то в случае, если обработка документа произошла с артефактом, то существует риск извлечения слова с разрывом <<Ref erences>>. Пробел тут ломает маску шаблона\cite{tellez2005machine}.

Сам список литературы может начинаться со множества вариантов: <<Список литературы>>, <<Библиография>>, <<Литература>>, References, Bibliography, Cited works, -- или такого заголовка вообще нет. Это затрудняет создание универсальной маски для поиска позиции библиографии и показывает, что априорно не известно какой вариант задействован в случайном журнале\cite{tkaczyk2015cermine}. Заголовок начала секции с библиографией может быть написан прописными буквами, с точкой в конце, в одну строку с горизонтальной чертой или вообще отсутствовать.

При попытке детекции позиции библиографии могут встречаться ложные срабатывания. Так, слово <<литература>> в статье может быть просто одним из чередующих слов в тексте, а может быть началом секции\cite{romary2015grobid}. 

Библиография не всегда идёт в конце статьи. Существуют варианты, когда в статьях с двухколонной вёрсткой колонтитулы вклиниваются между последним абзацы текста и заголовком списка литературы. Кроме того после списка литературы могу идти примечания редактора, благодарности или вовсе пустые страницы.

Помимо привычных способов оформления литературы в специальной для этого секции существует другой классический способ, а именно упоминания статьи в сносках. Такие сноски обычно располагаются в нижнем колонтитуле\cite{zhu2026benchmarking}.

\subsection{Проблемы извлечения библиографических записей}
После выявления позиции библиографии, оттуда необходимо извлечь библиографические записи. Проблема заключается в <<стилевом хаосе>>, артефактах парсинга и нестандартной структуры списка литературы\cite{ivanovic2016new}.

Неоднородность разделителей заключается в том, что разные журналы могут иметь разные способы отделения записей друг от друга, например, пустые строки, нумерованный список или маркированный, отсутствие явного разделителя, когда записи идут подряд. Если библиотека-парсер извлекла текст с нарушением структуры строк, то пустая строка может исчезнуть, а маркер списка прилипнуть к строке\cite{backes2024comparing}. 

Многострочные записи и висячие строки так же создают пробелы. Одна библиографическая запись может занимать несколько строк текста. При этом строчки с записями, которые начинаются с отступа, в извлечённом тексте могут передавать пробелами или табуляциями. С другой стороны эти отступы могу потеряться при парсинге документа или совпадать с началом записи, если нумерация смещена\cite{tkaczyk2015cermine}.

Артефакты, полученные с предыдущих этапов, как ошибки при извлечение текстового слоя или поиске области библиографии напрямую влияют на сегментацию. Так, если документ выполнен в двухколонной вёрстке, разрыв строки внутри записи может поделить на две колонки, и парсер соберёт их в неправильном порядке. Потеря номеров страниц или колонтитулов внутри области библиографии создаёт ложные пустые строки. Если область библиографии была обрезана слишком поздно, то первая или последняя запись будет неполной\cite{korner2017reference}.

Нестандартные способы оформления библиографии. Библиография может набрана в виде сплошного абзаца, где записи разделены точкой с запятой или специальным символом. Такой стиль может встречаться в старых изданиях или в журналах с жёсткой экономией места. Некоторые статьи используют сквозную нумерация ссылок по всему тексту, и с конце приводится список литературы под номерами, но без визуального разделения пустыми строками. Встречаются случаи, когда после библиографии идут примечания, благодарности или технические метки, ошибочно принимаемые за часть последней записи\cite{di2026tool}.

Из-за шумов OCR или некорректного распознавания текстового слоя две соседние записи могут слиться в одну. И наоборот, одна запись может разорваться на две из-за лишнего перевода строки или случайно попавшего колонтитула. Оба случая приводят к катастрофическом ошибкам на следующем этапе структурирования библиографических областей. 

\subsection{Проблемы структурирования библиографических областей}
Из каждой библиографической записи нужно выделить библиографические области. При работе с ними повторяется проблема <<стилевого хаоса>>. Одни и те же сущности, как автор, заголовок цитируемой статьи или источник могут иметь различное оформление. Области внутри записи могут быть неоднородны, следовать в разном порядке или иметь необязательные элементы, например, DOI или электронный адрес\cite{zhu2026benchmarking}.

Такая задача выделения библиографических областей из записи можно отнести к задаче парсинга ссылок (Reference Parsing или Citation Sequence Labeling). Её обычно сводят к задаче распознавания именованных сущностей (Named Entity Recognition)\cite{romary2015grobid}.

Некоторые библиографические области могут иметь вложенные или смешанные поля. Заголовок статьи может содержать знаки препинания, специальные символы, формулы, греческие буквы, например <<$\beta$-частицы>>. При парсинге записи такой текст может превратиться в последовательность токенов без семантики. Это затрудняет отделение заголовка упомянутой статьи от других областей. 

При работе с авторами можно столкнуться с неоднородностями: инициалы авторов, особенно если они не содержат знаки препинания, фамилии с частицами, например фамилии <<ван Дейк>> или <<де Брюйне>>, разделение авторов различной пунктуацией или союзами. Эти проблемы добавляют сложности при выделении блока автора.

Помимо всего этого артефакты предыдущих этапов вносят свой негативный вклад в этот этап\cite{tkaczyk2015cermine}. 
\begin{itemize}
    \item Шум, возникающий ещё на первом этапе, может сдвинуть границы библиографии, вследствие чего в на следующий этап могут попасть сторонние строки, ошибочно идентифицированные как библиографические записи, что приведет к формированию ложных областей.
    \item Если две записи слились из-за пропущенного разделителя, то поля второй записи ошибочно отнесут к первой.
    \item Разрыв строки внутри фамилии приводит к тому, что модель может не распознать корректно ни единый токен автора. 
\end{itemize}

\subsection{Проблемы галлюцинаторных списков}
В последнее время в конвейере разбора библиографии появилась новая проблема\cite{zhao2026llm}. Повсеместно применяют генеративный искусственный интеллект. В том числе в академическое письмо. Массовое применение БЯМ авторами научных публикацию привело к возникновению феномена генеративных галлюцинаций в списках литературы, что выражается в автоматическом создании правдоподобно выглядящих, но фактически фантомных библиографических записей. Модели с лёгкостью компилируют реальные фамилии учёных, валидные названия авторитетных журналов и типичные синтаксические конструкции стилей цитирования, создавая иллюзию академической обоснованности\cite{fiorillo2026confabulated}.

Для систем автоматического извлечения тенденция практик применения генеративного искусственного интеллекта меняет саму парадигму задачи. Извлечение структурных полей из галлюцинаторной ссылки будет содержать ложные идентификаторы, например несуществующие DOI, вымышленные номеры томов или несовпадающие годы изданий. Таким образом традиционный синтаксический анализ становиться недостаточным, порождая острую необходимость внедрения этапов внешней семантической верификации и перекрёстной сверки извлечённых токенов с глобальными индексами научного цитирования, например, Crossref, Scopus, PubMed\cite{abbonato2026checkifexist}.

Вот пример галлюцинаторного списка, созданного с помощью БЯМ, всех этих библиографических записей в реальности не существует:
\begin{enumerate}
    \item \textbf{Синтез авторов-экспертов (стиль APA):} \\
    Vaswani, A., Bengio, Y., \& Goodfellow, I. (2021). Generative Transformers for Sequence-to-Sequence Semantic Parsing. {\em IEEE Transactions on Pattern Analysis and Machine Intelligence}, 43(8), 2415--2429. DOI: 10.1109/TPAMI.2021.9876543. \\
    \textit{Пояснение:} Пример \guillemotleft химеры\guillemotright{}, где модель объединила реальных и крайне цитируемых ИИ-исследователей и поместила их вымышленную статью в авторитетный существующий журнал.

    \item \textbf{Имитация отечественных стандартов (стиль ГОСТ):} \\
    Кузнецов, С. И., Сидоров, П. А. Анализ структуры документов с использованием больших языковых моделей // {\em Автометрия}. --- 2024. --- Т. 60, № 3. --- С. 45--58. DOI: 10.15372/AUT20240305. \\
    \textit{Пояснение:} Классическая галлюцинация на русском языке. Модель полностью воспроизвела пунктуационную структуру ГОСТ и использовала реальный журнал СО РАН, но указанной статьи не существует.

    \item \textbf{Фантомное медицинское исследование:} \\
    Greenhalgh, T., \& Topol, E. J. (2023). The impact of large language models on systematic literature reviews in healthcare. {\em The Lancet Digital Health}, 5(6), e312--e320. DOI: 10.1016/S2589-7500(23)00123-4. \\
    \textit{Пояснение:} Модель использовала имена известных ученых из области цифровой медицины для генерации крайне правдоподобного названия статьи, релевантного профилю существующего журнала {\em The Lancet}.

    \item \textbf{Галлюцинация материалов конференции (стиль ACM):} \\
    Manning, C. D., Jurafsky, D., \& Ng, A. Y. (2022). Unsupervised Extraction of Bibliographic Metadata via Zero-Shot Prompting. In {\em Proceedings of the 60th Annual Meeting of the Association for Computational Linguistics (ACL 2022)}, 1402--1415. \\
    \textit{Пояснение:} Публикация в материалах известной NLP-конференции вымышлена ИИ-ассистентом на основе профильной терминологии авторов-экспертов.

    \item \textbf{Вымышленное монографическое издание:} \\
    Russell, S., \& Norvig, P. (2023). {\em Artificial Intelligence in Academic Writing: Semantic Web and Information Retrieval Perspectives}. Springer, Cham. DOI: 10.1007/978-3-031-12345-6. \\
    \textit{Пояснение:} Ссылка на несуществующую книгу авторов фундаментального учебника по ИИ в реальном научном издательстве Springer.
\end{enumerate}

Этот пример показывает, что в дальнейших работах в конвейер извлечения библиографии необходим блок проверки библиографических записей на существование.

\section{Обзор инструментов}
За историю разбора библиографии появилось и развивалось множество методов.

\subsection{Регулярные выражения}
Одними из первых методов являются регулярные выражения (Regular expressions, RegEx) -- специальный формальный язык поиска и манипуляции подстроками в тексте, основанный на использование метасимволов или масок, то есть можно сказать, что инструмент сопоставления текста с заданным шаблоном. Регулярные выражения являются хрупким решением, работающим в довольно жёстких ограниченных входных условиях, и при смене шаблона оформления способны перестать работать\cite{davis2019rethinking}.

Регулярные выражения нередко используют для определения оффсетов границ позиции библиографии или библиографических ссылок, чтобы потом передать полученную секцию для обработки более мощному инструменту. Библиографические области в библиографических записях в списках литературы обычно идут сплошным текстом. Регулярному выражению сложно понять, где заканчивается название статьи и начинается название журнала, если они разделены только запятой или точкой. 

При работе с регулярными выражениями нужно учитывать множество деталей. В зависимости от стиля оформления библиографии в журнале различные сущности могут оформляться по разному. Например, ФИО автора публикации может записываться несколькими способами: <<Красоткин С.А.>>, <<Красоткин, С.>>, <<Красоткин Семён>>, <<С.А. Красоткин>>. Довольно тяжело предусмотреть все инициалы и знаки препинания в одном правиле.

Необходимый учёт множества деталей в правилах регулярного выражения с условиям того, что различные издательства могут использовать свой собственный стиль: ГОСТ, APA, MLA, Harvard, Chicago. Так шаблон, заточенный под ГОСТ, сломается на стиле APA. Из-за чего почти под каждый стиль писать своё регулярное выражение.

Среди сильных сторон регулярных выражений выделяют: высокую скорость работы особенно при использовании детерминированных конечных автоматов, способные за счёт низкоуровневых алгоритмов обрабатывать большие массивы текстов за миллисекунды; кроссплатформеность, компактность шаблонов, заменяющая несколько строчков кода одним регулярным выражением. 

% Добавить расшифровку NFA, RE2
Высокая скорость работы регулярных выражений достигается на оптимизированных выражениях и простых задачах. Некоторая механизмы, например, NFA с использованием backtracking могут работать очень долго. Для борьбы с этим существуют библиотеки: RE2\cite{horky2022automata}, Hyperscan\cite{shuai2021performance}, обеспечивающие линейную скорость O(n), но они ради гарантированной линейной скорости и предсказуемости жертвуют некоторыми механизмами, например, обратными ссылками.

Компактность регулярных выражений не всегда является благом. Сложные якобы универсальные регулярные выражения плохо читаемы. Из-за чего их трудно читать, редактировать и поддерживать. Минимальные изменения, как лишний пробел или перенос строки, в структуре регулярного выражения могут полностью сломать работу шаблона. В дополнение можно подметить, что регулярные выражения не понимают смысл слов. Они ищут совпадение знаков, поэтому возможны ложноположительные срабатывания или, при попытке их избежать, к ложноотрицательным. Это происходит потому, что регулярные выражения сопоставляет строки на основе заданных форм, а не на основе семантического контекста.

\subsection{GROBID}

Одним из базовых инструментов извлечения библиографических данных остаётся инструмент на основе CRF (Conditional Random Fields) и BiLSTM-CRF с Glove эмбеддингами под названием GROBID (GeneRation Of Bibliographic Data)\cite{romary2015grobid}. Эта библиотека принимает на вход PDF-файл в структурированный TEI/XML.

GROBID демонстрирует высокую точность на англоязычных публикациях\cite{tkaczyk2018machine}, достигая значение F1-меры 0,89, но его качество существенно снижается при работе с:
\begin{itemize}
    \item Многоязычными документами;
    \item Ссылками в сносках (footnotes), хотя разработчики по этому вопросу пишут: <<поддерживаются, но всё ещё в разработке>>;
    \item Историческими стилями цитирования;
    \item Сложной вёрсткой, впрочем в работе\cite{meuschke2023benchmark} отмечается, что на данный момент почти все известные инструменты испытывают трудности с извлечением списков, колонтитулов и уравнений.  
\end{itemize}

Последние работы рекомендуют гибридное развёртывание\cite{zhu2026benchmarking}. Для хорошо структурированных PDF-файлов использовать GROBID, а многоязычные документы перенаправлять на вход БЯМ.

GROBID работает с текстовым слоем PDF и на документах-сканах не будет. К тому же для улучшения точности на сложных страницах нужны предварительные обработчики. Следовательно, для лучшей работы с GROBID лучше предварительно прогонять документы через OCR.

\subsubsection{Мультимодальные модели для анализа документов}

В ответ на ограничения текстовых подходов появились мультимодальные архи\-тек\-ту\-ры, объединяющие текстовые и визуальные признаки. 

Модели семейства LayoutLM (v2/v3) демонстрируют результаты на уровне SOTA на задачах анализа макета, однако имеют жесткие ограничения по лицензии для промышленного использования. Архитектура LayoutLMv3 представляет собой комбинацию текстового энкодера на базе RoBERTa и визуального энкодера ResNet. Входные данные включают три типа эмбеддингов: токен-эмбеддинги, позиционные эмбеддинги для координат bounding box'ов и эмбеддинг изображений.

Модель LiLT (Language-independent Layout Transformer) решает проблему языковой зависимости за счет разделения потоков обработки текста и макета, что позволяет эффективно применять трансферное обучение на многоязычных корпусах. Ключевое отличие LiLT заключается в том, что текстовая ветвь обучается на многоязычных данных независимо от визуальной. Это позволяет применять трансфер лёрнинг: визуальная компонента, обученная на англоязычных документах, успешно адаптируется к китайским или русским текстам без переобучения. 

Для задачи извлечения библиографии мультимодальные модели демонстрируют явные преимущества в сценариях со сложной версткой: двухколоночные статьи, врезки, таблицы внутри списка литературы, где чисто текстовые подходы терпят неудачу из-за нарушения порядка токенов. Однако их промышленное применение ограничено лицензиями. LayoutLMv3 распространяется под Microsoft Research License. Кроме того, данный класс решений крайне требователен к вычислительным ресурсам: дообучение на уровне документа требует не менее 16~ГБ видеопамяти (VRAM).

В работе~\cite{joshi2023end} предлагается сквозной (end-to-end) конвейер на основе LiLT для детекции библиографического раздела и fine-tuned модели SciBERT для сегментации отдельных ссылок. На внутреннем датасете Elsevier авторы достигли показателя $F_1 = 0{,}946$, превзойдя инструмент GROBID. К преимуществам данного подхода относятся полный контроль над пайплайном (в отличие от закрытых black-box решений), потенциал для расширения на многоязычные и сканированные документы, а также возможность модульной замены отдельных компонент без перестройки всей системы. Тем не менее, метод критически требователен к наличию больших объемов размеченных данных для обучения и остается чувствительным к качеству первичного PDF-парсинга.

\subsection{Большие Языковые Модели}
Для работы с библиографией используются также большие языковые модели. Например, семейства Qwen2.5. При подаче отдельных записей наблюдается  явление Zero-shot точность извлечения стремится к 100\%.

Работа \cite{zhu2026benchmarking} представляет систематическое сравнение современных БЯМ: DeepSeek-V3.1, Mistral-Small-3.2-24B, Gemma-3-27B-it, Qwen3-VL, -- с GROBID на трёх наборах данных, отражающих специфику социальных и гуманитарных наук: 
\begin{itemize}
    \item CEX: 112 англоязычных статей, 27 дисциплин;
    \item EXCITE: 351 документ на немецком/английском, включая ссылки в сносках и смешанные форматы;
    \item LinkedBooks: 1194 размечанных ссылки из гуманитарны публикаций с высокой стилистической и языковой вариативностью: $72,3\%$ итальянский, $18,3\%$ английский.
\end{itemize}

Извлечение ссылок насыщается при умеренной ёмкости модели: большинство БЯМ достигают сопоставимых результатов $F1 > 0,93$. Различия определяются скорее скоростью и частотой сбоев.

Парсинг ссылок и end-to-end обработка остаются узкими местами: здесь важна способность модели генерировать структурированный JSON-вывод в условиях шума и длинного контекста.

На наборах данных, ориентированных на социально-гуманитарный домен, как EXCITE или LinkedBooks адаптированные БЯМ часто превосходят GROBID, демонстрируя лучшую робастность к распределительным сдвигам.

При работе с большими языковыми моделями возникает проблема дефицита окна контекста при обработке полных текстов статей. Существует риск OOM. Это можно решать локальным анализом изолированных библиографических фрагментов.

Сегментация документов повышает надёжность end-to-end парсинга за счёт локализации ошибок и сокращения длины генерируемого вывода.

\subsection{Другие подходы}

В статье\cite{kopanichuk2025structure} описан гибридный метод сегментации текста, комбинирующий эвристические правила и градиентный бустинг. Он использует регулярные выражения для идентификации границы списка литературы. Не способен учесть все варианты оформления. Показывает низкое качество извлечения  авторов с частицами и короткими фамилиями. Формат API удобен для написания скрипта.

%CERIMNE, Neural ParsCit, AnyStyle.

\subsection{Сравнение подходов}

Для сис\-те\-ма\-ти\-за\-ции результатов проведенного обзора существующих ин\-стру\-мен\-тов це\-ле\-со\-об\-раз\-но провести их многокритериальный анализ. Сравнительная оценка различных технологических подходов осуществлялась на основе комплекса практических требований, критических для современных систем анализа научной деятельности: степени поддержки кириллического текста (в частности, русскоязычной специфики), способности обработки графических сканов документов, объема и качества размеченных данных, необходимых для дообучения, лицензионной чистоты для коммерческого использования, а также возможности автономного локального развертывания без обращения к внешним облачным интерфейсам (API).

Детальный сопоставительный анализ выявляет ряд фундаментальных за\-ко\-но\-мер\-нос\-тей, определяющих границы применимости каждого класса методов. Правило-ориентированные подходы, основанные на регулярных выражениях и детерминированной гибридной сегментации, демонстрируют низкий вычислительный порог входа и абсолютную предсказуемость результатов, однако их масштабирование ограничено необходимостью ручной экспертной адаптации под каждый новый стиль оформления библиографии. 

Статистические модели на основе условных случайных полей обеспечивают высокую точность извлечения в рамках своего целевого домена. Тем не менее, их адаптация к новым языкам или нестандартным стилям разметки требует повторного обучения с нуля, что становится трудновыполнимой задачей при отсутствии репрезентативных локальных корпусов. 

Большие языковые модели в режиме работы без предварительной настройки обеспечивают максимальную универсальность за счет колоссального масштаба предобучения. Однако их внедрение сопряжено с рисками генеративных галлюцинаций и высокими требованиями к аппаратным ресурсам. В свою очередь, мультимодальные архитектуры, теоретически наиболее полно утилизирующие пространственно-текстовую информацию, на практике ограничены жесткими лицензионными соглашениями и экстремальной сложностью развертывания на локальных вычислительных узлах.

В современной практике интеллектуального анализа текстов наметился отчетливый тренд к преодолению указанных компромиссов посредством внедрения динамической архитектуры на основе модели маршрутизатора запросов. В рамках такой парадигмы поток входящих документов дифференцируется в зависимости от их структурной сложности. Простые документы, англоязычные статьи со стандартной двухколоночной версткой, направляются в вычислительно легковесный и точный инструмент GROBID. Сложные же случаи, характеризующиеся наличием постраничных сносок, многоязычными метаданными или редкими историческими стилями цитирования, перенаправляются для обработки в адаптированные большие языковые модели. Именно такой гибридный подход позволяет достичь баланса между вычислительной эффективностью и полнотой извлечения данных.

\section{Дообучение Больших Языковых Моделей}

Большие языковые модели (БЯМ) демонстрируют выдающиеся способности к обобщению в Zero-shot режиме, однако их применение к специализированным задачам извлечения структурированных данных сопряжено с систематическими отклонениями. 

Формально задачу дообучения можно сформулировать следующим образом\cite{hu2022lora}. Пусть $\theta_0 \in \mathbb{R}^d$ --- параметры предобученной модели, где $d$ --- размерность пространства весов (для Llama-3-8B $d \approx 8 \times 10^9$). Дана обучающая выборка пар $\mathcal{D}_{\text{train}} = \{(x_i, y_i)\}_{i=1}^N$, где $x_i$ --- входная библиографическая запись, $y_i$ --- целевой JSON-объект с разметкой полей. Требуется найти адаптированные параметры $\theta^*$, минимизирующие функцию потерь на downstream-задаче:
\begin{equation}
\theta^* = \arg\min_{\theta} \sum_{(x,y) \in \mathcal{D}_{\text{train}}} \mathcal{L}\left(f_{\theta}(x), y\right) + \lambda \mathcal{R}(\theta, \theta_0),
\end{equation}
где $\mathcal{L}$ --- функция потерь (типично кросс-энтропия для задачи генерации), $\mathcal{R}$ --- регуляризатор, отражающий желаемую близость к предобученным весам, $\lambda$ --- коэффициент регуляризации.

Наиболее прямолинейный подход -- полное дообучение, при котором обновляются все параметры $\theta$. Однако для моделей семейства Llama-3 данный подход практически неосуществим по ряду причин \cite{liu2023parameter}.

Для модели с 8 миллиардами параметров в формате bfloat16 требуется 16~ГБ памяти только на хранение весов. Обучение с использованием метода обратного распространения ошибки требует дополнительного хранения градиентов (ещё 16~ГБ), состояний оптимизатора (Adam: до 32~ГБ) и активаций промежуточных слоёв (8--32~ГБ в зависимости от размера батча и длины контекста). Совокупное потребление VRAM без применения техник оптимизации памяти превышает 40~ГБ, что выходит за рамки массовых GPU-ускорителей (например, NVIDIA RTX~4090 располагает 24~ГБ) \cite{rajbhandari2019zero, liu2023parameter}.

Полное дообучение на узком корпусе библиографических записей (даже при $\lambda > 0$) приводит к смещению распределения весов $\theta$ в сторону целевого домена. Эмпирически установлено, что после FFT на доменных выборках модель демонстрирует статистически значимую деградацию в общих задачах понимания текста и генерации, включая падение метрик на бенчмарках типа HellaSwag и MMLU, что критично для гибридных систем, где LLM обрабатывает вспомогательные запросы \cite{hu2022lora, liu2023parameter}.

Каждая итерация дообучения порождает полную копию весов $\theta^*$ (16~ГБ в bf16), что усложняет версионирование, A/B-тестирование и откат к предыдущим версиям в production-среде. В отличие от методов эффективной адаптации, где сохраняется один базовый чекпоинт, а для каждой задачи хранятся лишь компактные дельта-модули \cite{hu2022lora}\cite{kaddour2023challenges}.

Альтернативой FFT служат методы параметрически эффективного дообучения (Parameter-Efficient Fine-Tuning, PEFT), при которых замораживается основная масса весов $\theta_0$ и обучается небольшое число дополнительных параметров $\Delta\theta$, где $|\Delta\theta| \ll |\theta_0|$ \cite{liu2023parameter}. К наиболее развитым подходам относятся:
\begin{itemize}
    \item Вместо модификации весов модели оптимизируются встраивания (embeddings) специальных токенов, добавляемых к входному промпту \cite{lester2021power}. Для задачи извлечения библиографии данный метод показал неудовлетворительные результаты: модель не способна освоить сложную структурную трансформацию <<сырая строка $\rightarrow$ JSON>> через манипуляцию лишь входным пространством.
    \item Обучаемые префиксы добавляются к ключам и значениям механизма внимания (attention) в каждом слое трансформера \cite{li2021prefix}. Метод эффективен для классификации текста, однако для генерации структурированного JSON с переменной длиной полей требует значительного числа префиксных параметров (1--2\% от $\theta_0$), что частично нивелирует преимущества PEFT.
    \item В архитектуру трансформера внедряются небольшие полносвязные слои (bottleneck-структура с понижением размерности) после каждого субслоя self-attention и feed-forward \cite{houlsby2019parameter}\cite{pfeiffer2020adapterhub}. Adapter Tuning демонстрирует хорошие результаты при $|\Delta\theta| \approx 2\text{--}4\%$ от $\theta_0$, однако вносит задержку инференса: каждый adapter требует дополнительных матричных умножений на этапе генерации, что нарушает конвейерную оптимизацию современных фреймворков вывода \cite{liu2023parameter}.
\end{itemize}

Наиболее перспективным для задачи извлечения библиографии признан метод низкоранговой адаптации LoRA (Low-Rank Adaptation), предложенный Hu et al.~(2021) \cite{hu2022lora}. Математическая суть метода заключается в следующем.

Пусть $W_0 \in \mathbb{R}^{d \times k}$ --- замороженная матрица весов одного из линейных слоёв предобученной модели (например, проекция query в механизме внимания). Вместо прямого обучения приращения $\Delta W$ (что эквивалентно FFT), LoRA параметризует $\Delta W$ через произведение двух матриц низкого ранга:
\begin{equation}
W = W_0 + \Delta W = W_0 + \frac{\alpha}{r} \cdot B \cdot A,
\end{equation}
где $B \in \mathbb{R}^{d \times r}$, $A \in \mathbb{R}^{r \times k}$, $r \ll \min(d, k)$ --- гиперпараметр ранга адаптера, $\alpha \in \mathbb{R}$ --- масштабирующий коэффициент. При инициализации $A$ заполняется случайными значениями из нормального распределения $\mathcal{N}(0, 1)$, а $B$ --- нулями, что обеспечивает $\Delta W = 0$ на старте обучения и точное сохранение поведения базовой модели \cite{hu2022lora}.

Для Llama-3-8B с $d = 4096$ (размерность скрытого состояния) и стандартным рангом $r = 16$ количество обучаемых параметров на один линейный слой (при условии $k=d$ для проекционных матриц) составляет $d \cdot r + r \cdot k = 2 \cdot 4096 \cdot 16 = 131072$ вместо $4096 \cdot 4096 = 16777216$, то есть сжатие в 128 раз. Совокупно для всех целевых модулей механизма внимания при $r = 16$ обучается $\approx 16{,}8$ миллиона параметров --- лишь 0{,}2\% от 8 миллиардов базовой модели \cite{hu2022lora, liu2023parameter}.

На этапе эксплуатации матрицы $B \cdot A$ могут быть предварительно вычислены и слиты с $W_0$ (merge adapters): $W_{\text{merged}} = W_0 + \frac{\alpha}{r} B A$. Таким образом, структура вычислений идентична базовой модели, что критично для production-среды ИСАНД с требованиями к latency \cite{hu2022lora}\cite{meng2023efficient}.

Различные LoRA-адаптеры могут храниться как независимые модули небольшого размера ($\approx$~30~МБ для $r=16$ против 16~ГБ полной модели) и динамически подгружаться для разных задач: адаптер для русскоязычных записей по ГОСТ, адаптер для англоязычных записей APA, адаптер для смешанных корпусов. Данное свойство обеспечивает масштабируемость системы без дублирования базовых весов \cite{hu2022lora}\cite{wang2023lorahub}.

Для преодоления ограничений по памяти на этапе обучения применяется модификация QLoRA (Quantized LoRA), предложенная Dettmers et al.~(2023) \cite{dettmers2023qlora}. Метод комбинирует три ключевые техники:
\begin{itemize}
    \item 4-битное квантование базовых весов. Матрица $W_0$ хранится в сжатом формате NF4 (Normal Float 4), где каждый вес кодируется 4 битами вместо 16 (bfloat16) или 32 (float32). Динамический диапазон квантоватора адаптируется под эмпирическое распределение весов трансформера, что минимизирует ошибку квантования \cite{dettmers2023qlora}.
    \item Двойная квантизация. Параметры квантователя (scaling factors) дополнительно сжимаются с 32 до 8 бит, что сокращает накладные расходы на хранение метаданных \cite{dettmers2023qlora}.
    \item Страничный оптимизатор. Градиенты, состояния оптимизатора (для AdamW --- первый и второй моменты) и активации хранятся в pageable unified memory с автоматической выгрузкой на CPU при нехватке VRAM, что предотвращает out-of-memory ошибки без ручного управления памятью \cite{dettmers2023qlora, rajbhandari2019zero}.
\end{itemize}

Для Llama-3-8B полное дообучение требует $\approx 48$~ГБ VRAM (при батче 1 и длине контекста 512). QLoRA с 4-битным квантованием, рангом $r=16$ и батчем 2 умещается в 10{,}5~ГБ VRAM, что позволяет проводить обучение на GPU-ускорителе NVIDIA Tesla~T4 (16~ГБ) --- стандартном облачном инстансе. Скорость обучения снижается на 35--40\% по сравнению с полной точностью из-за накладных расходов на де-квантование в прямом и обратном проходах, однако абсолютное время обучения на корпусе из 20~000 записей (100 шагов, батч 2, накопление градиентов на 4 шага) составляет $\approx 2{,}5$ часа, что приемлемо для итеративной разработки \cite{dettmers2023qlora, liu2023parameter}.

Адаптация LoRA/QLoRA к задаче структурирования библиографических записей обладает специфическими особенностями, вытекающими из природы целевых данных.

В типичной библиографической записи поле <<Авторы>> занимает 15--25\% токенов, <<Название статьи>> -- 20--40\%, тогда как <<Год>>, <<Том>>, <<Номер>>, <<Страницы>> -- каждое 2--5\%. Стандартная кросс-энтропия приводит к доминированию градиентов от частых классов. Для компенсации применяется взвешивание классов обратно пропорционально их частоте в обучающей выборке \cite{chawla2002smote, jiang2023llm_imbalance}:
\begin{equation}
w_c = \frac{N_{\text{total}}}{N_c \cdot |\mathcal{C}|},
\end{equation}
где $N_c$ --- число токенов класса $c$, $|\mathcal{C}| = 7$ -- число библиографических областей.

В отличие от свободной генерации текста, библиографическая запись должна соответствовать жёсткой JSON-схеме с обязательными ключами. Для обеспечения валидности на этапе обучения применяется техника constrained decoding: логиты модели маскируются так, что на каждом шаге генерации допустим только токен, сохраняющий синтаксическую корректность частичного JSON (открывающие/закрывающие скобки, кавычки, запятые, допустимые литералы \texttt{null}). Эмпирически установлено, что данный подход снижает долю синтаксических ошибок более чем на 90\% по сравнению с неограниченной генерацией \cite{willard2023efficient, geng2023constrained}.

Дефекты входного датасета (дублирование системного промпта в поле \texttt{output}) приводят к катастрофической деградации: модель обучается циклически воспроизводить инструкцию вместо анализа. Для предотвращения применяется многоуровневая валидация данных: (1) синтаксическая проверка JSON, (2) семантическая проверка типов полей (год -- 4 цифры, страницы -- диапазон с тире), (3) кросс-валидация с эталонной разметкой на уровне 5\% hold-out выборки перед запуском обучения \cite{liu2023parameter, wang2024llm_data_quality}.

Стандартная конфигурация LoRA применяет адаптеры ко всем линейным слоям внимания и MLP-блоков. Для задачи извлечения структурированных данных экспериментально установлено, что наибольший прирост качества даёт адаптация проекций value ($v\_proj$) и output ($o\_proj$) в слоях внимания, тогда как адаптация MLP-блоков ($gate\_proj$, $up\_proj$, $down\_proj$) вносит вклад в 3--5\% Macro $F_1$ при удвоении числа обучаемых параметров. Оптимальная конфигурация для данной задачи: все линейные слои с рангом $r=16$, что обеспечивает баланс между качеством и вычислительными затратами \cite{hu2022lora, kumar2024lora_target}.